\section{Model Dwukryterialny}
\paragraph{Zbiory}
\begin{itemize}
    \item $i, k \in F$ - Fabryki 
    \item $i, l \in M$ - Magazyny 
    \item $j \in K$ - Klienci
\end{itemize}
\paragraph{Parametery}
\begin{itemize}
    \item $C_{i, j}$ - Koszty transportu dóbr z punktów $i$ (fabryka lub magazyn) do klienta $j$
    \item $C_{k, l}$ - Koszty transportu dóbr z fabryki $i$ do magazynu $k$ 
    \item $P_{i, j}$ - Binarnie określa czy preferencja klienta zostąła spełniona (1) czy nie (0)
    \item $S_j$ - Poziom satysfakcji klienta $j$ - poziom satysfakcji liczymy w zależności od tego ile dany klient zamówił towaru \\ 
    Zadowolenie klientów którzy zamawiają więcej towarów jest traktowane priorytetowo:
    \[ S_1 = \frac{50}{60}  \]
    \[ S_2 = \frac{10}{60}  \]
    \[ S_3 = \frac{40}{60}  \]
    \[ S_4 = \frac{35}{60}  \]
    \[ S_5 = \frac{60}{60}  \]
    \[ S_6 = \frac{20}{60}  \]
    \item $D_j$ - Zapotrzebowanie klienta $j$
\end{itemize}
\paragraph{Zmienne decyzyjne}
\begin{itemize}
    \item $x_{i,j}$ - Liczba dóbr (w tys. ton) przetransportowana z punktu $i$ (fabryka lub magazyn) do klienta $j$
    \item $y_{k, l}$ - Liczba dóbr (w tys. ton) przetransportowana z fabryki $k$ do magazynu $l$
\end{itemize} 
\paragraph{Funkcja celu}
\begin{enumerate}
    \item Minimalizacja kosztów dystrybucji \\ 
    \[ Min(\sum_{i, j} C_{i, j} * x_{i, j} + \sum{i, k} C_{i, k} * y_{i, k}) \]
    \item Maksymalizacja satysfakcja klienta \\ 
    W celu maksymalizacji satysfakcji klienta policzymy ile z dostaw do klientów odbyło się z preferencyjnych źródeł \\
    \[ Max(\sum_{j} S_j * P_{i,j} * x_{i, j}) \] 
\end{enumerate}
\paragraph{Funkcja celu alpha beta}
\[ \alpha * (Min(\sum_{i, j} C_{i, j} * x_{i, j} + \sum{i, k} C_{i, k} * y_{i, k})) + \beta * (Max(\sum_{j} S_j * P_{i,j} * x_{i, j}))  \]
\paragraph{Ograniczenia}
Miesięczne możliwości produkcyjne fabryk
\begin{equation}
    \sum_j x_{F1, j} + \sum_k y_{F1, k} \leq 150
\end{equation}
\begin{equation}
    \sum_j x_{F2, j} + \sum_k y_{F2, k} \leq 200
\end{equation}
Miesięczna ilośc obsługiwanego towaru przez magazyny 
\begin{equation}
    \sum_j x_{M1, j} \leq 70
\end{equation}
\begin{equation}
    \sum_j x_{M2, j} \leq 50
\end{equation}
\begin{equation}
    \sum_j x_{M3, j} \leq 100
\end{equation}
\begin{equation}
    \sum_j x_{M4, j} \leq 40
\end{equation}
Spełnienie preferencji klienta 
\begin{equation}
    \sum_i x_{i, j} = D_j
\end{equation}
Wartości niezerowe 
\begin{equation}
    x_{i, j}, y_{i, k} \geq 0
\end{equation}

